\chapter{Literature Review}

\section{Introduction to the Field}
[This section should provide an overview of the research field and establish the context for your work. Include key concepts and definitions.]

\section{Historical Development}
The field of [your research area] has evolved significantly over the past several decades. Early work by \textcite{smith2020fundamentals} established the fundamental principles that continue to influence current research directions. Subsequent studies by \textcite{brown2019advances} expanded our understanding of the underlying mechanisms.

\section{Current State of Knowledge}
Recent advances in [your field] have been driven by several key developments:

\subsection{Theoretical Foundations}
The theoretical framework for understanding [your topic] was established through the work of \textcite{wilson2021theoretical}, who provided a comprehensive model that explains the fundamental mechanisms involved.

\subsection{Experimental Methods}
Experimental approaches have evolved significantly, with \textcite{anderson2022experimental} introducing novel techniques that have improved our ability to measure and characterize the phenomena under investigation.

\section{Key Research Areas}
Several important areas of research have emerged in recent years:

\subsection{Area 1: [Specific Research Area]}
[Description of the first research area with relevant citations and examples.] Recent work by \textcite{thompson2020comprehensive} has provided valuable insights into this field.

\subsection{Area 2: [Another Research Area]}
[Description of the second research area with relevant citations and examples.] \textcite{rodriguez2021novel} has contributed significantly to our understanding of this area.

\section{Research Gaps and Opportunities}
Despite significant progress, several important gaps remain in our understanding:

\begin{enumerate}
\item Gap 1: Describe what is not yet understood or investigated
\item Gap 2: Another area that needs further research
\item Gap 3: Additional research opportunities
\end{enumerate}

\section{Example Figures}
This section demonstrates how to include figures in your literature review. Figure \ref{fig:literature_overview} provides an overview of the research landscape, while Figure \ref{fig:research_trends} shows current trends in the field.

\begin{figure}[h!tbp]
\centering
\includegraphics[width=0.7\textwidth]{chapters/2/figures/fig_1.png}
\caption{Literature overview and research landscape. This figure demonstrates how to include and reference figures in your literature review chapter.}
\label{fig:literature_overview}
\end{figure}

\begin{figure}[h!tbp]
\centering
\includegraphics[width=0.7\textwidth]{chapters/2/figures/fig_2.png}
\caption{Current research trends and developments. This shows the evolution of research in the field over time.}
\label{fig:research_trends}
\end{figure}

\section{Summary}
[Provide a summary of the key points from the literature review and how they relate to your research objectives. This should be 1-2 paragraphs.]